% Source: https://artofproblemsolving.com/community/c3h2011108_a_geometry_problem
% Stable identity: aops-topic-2011108-q01
\textbf{Congruent triangles.} Compare triangles $ABE$ and $CAD$:
\[AB=CA,\qquad \angle ABE=\angle CAD=60^\circ,\qquad BE=AD.\]
By SAS they are congruent, with $A\leftrightarrow C$, $B\leftrightarrow A$, $E\leftrightarrow D$. In particular $\angle BAE=\angle ACD$.

\begin{asy}
usepackage("amsmath");
size(5.5cm);
defaultpen(fontsize(10pt));
pair A=(0,0), B=(6,0), C=rotate(60)*(6,0);
real t=2.2;
pair D=(t,0);
pair Ept=B+t*unit(C-B);
pair F=extension(A,Ept,C,D);
draw(A--B--C--cycle,linewidth(0.8));
draw(A--Ept,blue); draw(C--D,red);
void tick(pair P, pair Q){pair m=(P+Q)/2, u=unit(Q-P)*0.18; draw(m-rotate(90)*u--m+rotate(90)*u,linewidth(0.8));}
tick(A,D); tick(B,Ept);
draw(arc(F,0.55,degrees(Ept-F),degrees(C-F)),linewidth(0.7));
dot(A);dot(B);dot(C);dot(D);dot(Ept);dot(F);
label("$A$",A,SW); label("$B$",B,SE); label("$C$",C,N);
label("$D$",D,S); label("$E$",Ept,NE); label("$F$",F,W);
label("$60^\circ$",F+0.95*dir((degrees(Ept-F)+degrees(C-F))/2),fontsize(8pt));
label("\mbox{}",(min(currentpicture,user=true)+max(currentpicture,user=true))/2,black); // keeps the text after the figure black when LaTeX typesets the labels inline
\end{asy}

\textbf{Exterior angle.} The cevians $AE$ and $CD$ meet at an interior point $F$, which lies between $A$ and $E$. So $\angle CFE$ is the exterior angle at $F$ of triangle $AFC$, and
\[\angle CFE=\angle FAC+\angle FCA=\angle EAC+\angle ACD.\]
Replacing $\angle ACD$ by the equal angle $\angle BAE$,
\[\angle CFE=\angle EAC+\angle BAE=\angle BAC=60^\circ.\]
Hence $\angle CFE=\boxed{60^\circ}$, whatever the common length $AD=BE$.
